\documentclass[11pt, a4paper]{article}

% --- PREAMBLE ---
\usepackage[a4paper, top=2.5cm, bottom=2.5cm, left=2cm, right=2cm]{geometry}
\usepackage{fontspec}
\usepackage[english]{babel}

% Fonts setup for LuaLaTeX
\setmainfont{Noto Sans}
\setsansfont{Noto Sans}
\setmonofont{Noto Sans Mono}

% Packages for layout and formatting
\usepackage{graphicx}     % For images
\usepackage{titlesec}     % For section formatting
\usepackage{fancyhdr}     % For headers/footers
\usepackage{tabularx}     % For flexible tables
\usepackage{longtable}    % For tables that span pages
\usepackage{booktabs}     % For professional table lines
\usepackage{float}        % For figure placement
\usepackage{hyperref}     % For clickable links in TOC
\usepackage{enumitem}     % For better list control
\usepackage{titlesec}

% Hyperlink setup
\hypersetup{
    colorlinks=true,
    linkcolor=black,
    filecolor=magenta,      
    urlcolor=blue,
    pdftitle={Software Design Specifications - Orbis AI},
    pdfpagemode=FullScreen,
}


\titleformat{\paragraph}
  {\normalfont\normalsize\bfseries}
  {\theparagraph}{1em}{}
\titlespacing*{\paragraph}{0pt}{1em}{0.5em}

% Header and Footer setup
\pagestyle{fancy}
\fancyhf{}
\lhead{Software Design Specifications}
\rhead{Version: 1.5}
\cfoot{\thepage}

% Macros for project details
\newcommand{\projectname}{Orbis AI}
\newcommand{\projectversion}{1.5}
\newcommand{\supervisor}{Ms. Bakhtawar Abbasi}

% Macro for Data Dictionary Table
% Arguments: {TableName}{TableContent}
\newcommand{\datadictionarytable}[2]{
    \noindent\textbf{Table Name: #1}
    \begin{center}
    \begin{tabularx}{\linewidth}{|l|l|l|X|}
        \hline
        \textbf{Attribute} & \textbf{Data Type} & \textbf{Constraints} & \textbf{Description} \\
        \hline
        #2
    \end{tabularx}
    \end{center}
    \vspace{0.5cm}
}

% Macro for Diagram Placeholder
\newcommand{\diagramplaceholder}[2]{
    \begin{figure}[H]
        \centering
        \framebox{\parbox{0.8\linewidth}{\centering
            \vspace{4cm}
            \textbf{#1} \\
            \small\textit{#2}
            \vspace{4cm}
        }}
        \caption{#1}
    \end{figure}
}

\begin{document}

% ==========================================
% TITLE PAGE
% ==========================================
\begin{titlepage}
    \centering
    \vspace*{2cm}
    
    {\Huge \textbf{Software Design Specifications}}\\[1cm]
    {\huge \textbf{\projectname}}\\[0.5cm]
    {\Large \textbf{Intelligent Travel Planning System}}\\[0.5cm]
    {\Large \textbf{Version: \projectversion}}\\[2cm]
    
    \begin{tabularx}{0.8\textwidth}{l X}
        \textbf{Project ID:} & F25-54 \\ [10pt]
        \textbf{Supervisor:} & \supervisor \\ [10pt]
        \textbf{Co Supervisor:} & - \\ [10pt]
        \textbf{Project Team:} & M. Talha Yousif (Leader) \\
         & Salik Ahmed Bhatti \\
         & Toheed Ali \\ [10pt]
        \textbf{Submission Date:} & November 30, 2025 \\
    \end{tabularx}
    
    \vfill
    
    {\large \textbf{Department of Computer Science}}
    
    \vspace*{1cm}
    {\large \textbf{National University of Computer \& Emerging Sciences}}
\end{titlepage}

% ==========================================
% INSTRUCTIONS

% ==========================================
% DOCUMENT HISTORY
% ==========================================
\newpage
\section*{Document History}
\noindent [Revision history will be maintained to keep a track of changes done by anyone in the document.]

\vspace{0.5cm}

\begin{tabularx}{\textwidth}{|l|l|l|X|}
    \hline
    \textbf{Version} & \textbf{Author} & \textbf{Date} & \textbf{Description of change} \\
    \hline
    1.0 & M. Talha Yousif & 29/11/2025 & Initial document creation and structure setup \\
    \hline
    1.1 & M. Talha Yousif & 29/11/2025 & Introduction section completed with objectives \\
    \hline
    1.2 & Salik Ahmed Bhatti & 29/11/2025 & Project scope and exclusions defined \\
    \hline
    1.3 & Toheed Ali & 29/11/2025 & System constraints and interface requirements \\
    \hline
    1.4 & M. Talha Yousif & 30/11/2025 & Use case diagram and functional requirements \\
    \hline
    1.5 & Team & 30/11/2025 & Final review and SDS compilation \\
    \hline
\end{tabularx}

\newpage
\tableofcontents
\newpage

% ==========================================
% 1. INTRODUCTION
% ==========================================
\section{Introduction}

\subsection{Purpose of Document}
The purpose of this Software Design Specification (SDS) is to describe the system architecture, detailed design, and database structure of \textbf{Orbis AI}, an intelligent travel planning system. This document serves as a guide for developers, project managers, and quality assurance teams to understand the technical implementation of the multi-agent AI system and its integration with external travel APIs.

\subsection{Intended Audience}
\begin{itemize}
    \item \textbf{Developers:} To implement the backend logic, AI agents, and frontend interfaces.
    \item \textbf{Project Supervisors:} To evaluate the technical depth and design soundness.
    \item \textbf{QA Team:} To understand system flows for designing test cases.
\end{itemize}

\subsection{Document Convention}
\begin{itemize}
    \item \textbf{Database:} PostgreSQL conventions (snake\_case for tables and columns).
    \item \textbf{API:} RESTful standards for endpoints.
    \item \textbf{Variable Naming:} camelCase for JavaScript/React, snake\_case for Python backend.
\end{itemize}

\subsection{Project Overview}
Orbis AI is a comprehensive travel planning application that leverages a Multi-Agent System (MAS) to automate the complexities of trip organization. Unlike traditional platforms that require manual aggregation of flights, hotels, and activities, Orbis AI utilizes specialized AI agents (Coordinator, Flight, Hotel, Activity) to generate personalized itineraries. The system integrates with major travel APIs (Amadeus, Skyscanner) to provide real-time booking capabilities and dynamic itinerary adjustments.

\subsection{Scope}
The scope of Orbis AI includes:
\begin{itemize}
    \item \textbf{User Management:} Authentication, profile management, and travel history tracking.
    \item \textbf{AI Trip Planning:} Natural language processing for intent recognition and multi-agent coordination for itinerary generation.
    \item \textbf{Search \& Booking:} Real-time searching of flights and hotels via Amadeus/Skyscanner APIs and secure payment processing via Stripe.
    \item \textbf{Itinerary Management:} Interactive timelines, collaborative editing, and PDF export.
\end{itemize}
\textit{Exclusions:} Visa processing services, physical tour guide allocation, and offline-only functionality.

% ==========================================
% 2. DESIGN CONSIDERATIONS
% ==========================================
\section{Design Considerations}

\subsection{Assumptions and Dependencies}
\begin{itemize}
    \item \textbf{External APIs:} The system relies on the continuous availability of Amadeus, Skyscanner, and OpenAI/LangChain APIs.
    \item \textbf{Connectivity:} Users must have a stable internet connection to interact with the AI agents.
    \item \textbf{Payment Gateway:} Stripe test mode will be used for development; valid credentials are assumed for production.
\end{itemize}

\subsection{Risks and Volatile Areas}
\begin{itemize}
    \item \textbf{API Rate Limits:} Third-party travel APIs may impose rate limits that could affect search performance.
    \item \textbf{LLM Hallucinations:} AI agents might generate non-existent travel options; strict validation layers are required.
    \item \textbf{Data Privacy:} Handling sensitive user data (passports, payments) requires robust encryption and GDPR compliance.
\end{itemize}

\newpage
% ==========================================
% 3. SYSTEM ARCHITECTURE
% ==========================================
\section{System Architecture}

\subsection{System Level Architecture}
The system follows a modern client-server architecture with a specialized AI layer. The frontend is built with React/Next.js, communicating with a Python-based backend (FastAPI/Django) that orchestrates LangChain agents.
\begin{figure}[H]
    \centering
    \includegraphics[width=0.9\linewidth]{images/Architecture.png}
\end{figure}

\newpage
\subsection{Software Architecture}
Orbis AI utilizes a \textbf{Service-Oriented Architecture (SOA)} combined with an \textbf{Agentic Pattern}.
\begin{itemize}
    \item \textbf{Presentation Layer:} React.js web interface.
    \item \textbf{Application Layer:} Handles business logic, user sessions, and API orchestration.
    \item \textbf{Intelligence Layer:} Powered by LangChain and AutoGen, housing the distinct agents (flight search, hotel booking, itinerary optimizer).
    \item \textbf{Data Layer:} PostgreSQL with `pgvector` for vector embeddings (for personalized recommendations).
\end{itemize}

% ==========================================
% 4. DESIGN STRATEGY
% ==========================================
\section{Design Strategy}
The design strategy prioritizes \textbf{Modularity} and \textbf{Scalability}.
\begin{itemize}
    \item \textbf{Object-Oriented Design (OOD):} Used for defining agent classes and data models.
    \item \textbf{Agent-Based Modeling:} Decoupling the "Flight Search" logic from "Hotel Search" logic allows independent scaling and maintenance of agents.
    \item \textbf{Secure by Design:} Implementing Role-Based Access Control (RBAC) and data encryption at rest.
\end{itemize}

% ==========================================
% 5. DETAILED SYSTEM DESIGN
% ==========================================
\section{Detailed System Design}

\subsection{Database Design}

\subsubsection{ER Diagram}
\begin{figure}[H]
    \centering
    \includegraphics[width=0.9\linewidth]{images/ERD.png}
\end{figure}

\newpage
\subsubsection{Data Dictionary}
The database is implemented in PostgreSQL. Below are the key schemas derived from the project's SQL definition.

\paragraph{5.1.2.1 Users Table}
\datadictionarytable{public.users}{
    id & uuid & PK, Not Null & Unique user identifier \\ \hline
    email & text & Unique, Not Null & User email address \\ \hline
    full\_name & text & Nullable & User's full name \\ \hline
    status & user\_status & Default 'active' & Account status \\ \hline
    metadata & jsonb & Default '\{\}' & JSON store for arbitrary user data \\ \hline
    created\_at & timestamp & Default now() & Account creation time \\ \hline
}

\paragraph{5.1.2.2 Trips Table}
\datadictionarytable{public.trips}{
    id & uuid & PK, Not Null & Unique trip identifier \\ \hline
    user\_id & uuid & FK (users) & Owner of the trip \\ \hline
    title & text & Not Null & Name of the trip \\ \hline
    status & trip\_status & Default 'planning' & e.g., planning, booked, completed \\ \hline
    budget\_total & numeric & Check >= 0 & Total budget for the trip \\ \hline
    generated\_itinerary & jsonb & Nullable & The AI-generated itinerary structure \\ \hline
    is\_public & boolean & Default false & Visibility of the trip \\ \hline
}

\paragraph{5.1.2.3 Chat Sessions}
\datadictionarytable{public.chat\_sessions}{
    id & uuid & PK, Not Null & Unique session ID \\ \hline
    user\_id & uuid & FK (users) & User associated with the chat \\ \hline
    trip\_id & uuid & FK (trips), Nullable & Linked trip context \\ \hline
    title & text & Nullable & Summary or title of the chat \\ \hline
    context\_summary & text & Nullable & AI-generated summary of context \\ \hline
}

\paragraph{5.1.2.4 Messages}
\datadictionarytable{public.messages}{
    id & uuid & PK, Not Null & Unique message ID \\ \hline
    session\_id & uuid & FK (chat\_sessions) & Parent session \\ \hline
    sender\_type & text & Not Null & 'user', 'agent', or 'system' \\ \hline
    content & text & Not Null & The message body \\ \hline
    tool\_calls & jsonb & Default '[]' & Records of AI tool usage \\ \hline
}

\paragraph{5.1.2.5 Agent Runs}
\datadictionarytable{public.agent\_runs}{
    id & uuid & PK, Not Null & Unique execution ID \\ \hline
    session\_id & uuid & FK (chat\_sessions) & Context session \\ \hline
    agent\_type & text & Not Null & e.g., 'coordinator', 'flight\_agent' \\ \hline
    status & text & Not Null & 'running', 'completed', 'failed' \\ \hline
    tokens\_used & integer & Default 0 & Token consumption count \\ \hline
}

\paragraph{5.1.2.6 Flight Search Results}
\datadictionarytable{public.flight\_search\_results}{
    id & uuid & PK, Not Null & Unique result ID \\ \hline
    session\_id & uuid & FK & Associated chat session \\ \hline
    carrier & text & Not Null & Airline name \\ \hline
    departure\_time & timestamp & Not Null & Flight departure \\ \hline
    price & numeric & Not Null & Cost of the flight \\ \hline
    raw\_data & jsonb & Not Null & Full JSON response from API \\ \hline
}

\paragraph{5.1.2.7 Bookings}
\datadictionarytable{public.bookings}{
    id & uuid & PK, Not Null & Unique booking ID \\ \hline
    trip\_id & uuid & FK (trips) & Associated trip \\ \hline
    type & booking\_type & Not Null & 'flight', 'hotel', 'activity' \\ \hline
    provider\_booking\_id & text & Nullable & External ID from provider \\ \hline
    status & booking\_status & Default 'pending' & Status of the reservation \\ \hline
    cost & numeric & Not Null & Final cost \\ \hline
}

\paragraph{5.1.2.8 Payments}
\datadictionarytable{public.payments}{
    id & uuid & PK, Not Null & Unique payment ID \\ \hline
    booking\_id & uuid & FK (bookings) & Associated booking \\ \hline
    amount & numeric & Not Null & Payment amount \\ \hline
    currency & text & Default 'USD' & Payment currency \\ \hline
    stripe\_payment\_id & text & Unique & External Stripe Transaction ID \\ \hline
    status & payment\_status & Not Null & 'succeeded', 'failed', 'pending' \\ \hline
}

\paragraph{5.1.2.9 User Preferences}
\datadictionarytable{public.user\_preferences}{
    user\_id & uuid & PK, FK (users) & Linked user \\ \hline
    preferred\_airlines & text[] & Default '\{\}' & Array of airline codes \\ \hline
    dietary\_restrictions & text[] & Default '\{\}' & e.g., 'vegan', 'halal' \\ \hline
    accommodation\_type & text[] & Default '\{\}' & e.g., 'hotel', 'hostel' \\ \hline
    budget\_range & jsonb & Nullable & Min/Max preferences \\ \hline
}

\newpage
\subsection{Application Design}
\subsubsection{Sequence Diagrams}\vspace{0.3em}
\paragraph{5.2.1.1 User Onboarding \& Preference Setup}

\begin{figure}[H]
    \centering
    \includegraphics[width=0.9\linewidth]{images/User Onboarding.png}
\end{figure}

\paragraph{5.2.1.2 AI Trip Planning Workflow}
\begin{figure}[H]
    \centering
    \includegraphics[width=0.9\linewidth]{images/AI Trip Planning.png}
\end{figure}

\newpage
\paragraph{5.2.1.3 Booking \& Payment Process}
\begin{figure}[H]
    \centering
    \includegraphics[width=0.9\linewidth]{images/Booking & Payment.png}
\end{figure}

\newpage
\subsubsection{State Diagrams}

\paragraph{5.2.2.1 Trip Lifecycle State}
\begin{figure}[H]
    \centering
    \includegraphics[width=0.9\linewidth]{images/Trip Lifecycle.png}

\end{figure}

\newpage
\paragraph{5.2.2.2 Agent Execution State}
\begin{figure}[H]
    \centering
    \includegraphics[width=0.9\linewidth]{images/Agent Execution.png}
\end{figure}

\newpage
\paragraph{5.2.2.3 Booking Status State}
\begin{figure}[H]
    \centering
    \includegraphics[width=0.9\linewidth]{images/Booking Status.png}
\end{figure}


% ==========================================
% 6. REFERENCES
% ==========================================
\newpage
\section{References}
\begin{thebibliography}{9}

\bibitem{amadeus}
Amadeus for Developers. "Travel APIs." [Online]. Available: \url{https://developers.amadeus.com}

\bibitem{skyscanner}
Skyscanner. "Flights API." [Online]. Available: \url{https://skyscanner.net/developers}

\bibitem{booking}
Booking.com. "Distribution API." [Online]. Available: \url{https://developers.booking.com}

\bibitem{stripe}
Stripe. "Payment API (Test Mode)." [Online]. Available: \url{https://stripe.com/docs/api}

\bibitem{langchain}
LangChain. "Building Applications with LLM Agents." [Online]. Available: \url{https://python.langchain.com}

\bibitem{autogen}
Microsoft Research. "AutoGen: Enabling Next-Gen Multi-Agent AI Systems." In \textit{Proceedings of NeurIPS Workshops}, 2023.

\end{thebibliography}

% ==========================================
% 7. APPENDICES
% ==========================================
\section{Appendices}
\noindent This section is not applicable for the current document. No supplementary appendices are included.



\end{document}